\section{敏感性分析}
\label{sec:sens}

\subsection{返航安全余量对最大安全载荷与组批的影响}

返航电量下限（安全余量比例 $\eta_g$）是唯一直接由安全策略给定、且贯穿四问的关键参数。以附件基准 $\eta_g=20\%$ 为中心，在 $10\%\!\sim\!35\%$ 范围内扰动，重算最大安全载荷与组批方案。

对最大安全载荷 $q^{\max}$：机型 A 全程质量受限，$q^{\max}$ 恒为 $25\,$kg，对 $\eta$ 完全不敏感；机型 B 仅 $S_{008}$ 为能量受限，随 $\eta$ 增大略降但很快被质量上限 $30\,$kg 覆盖；机型 C 最敏感——其 $5$ 个能量受限服务区的临界载荷随 $\eta$ 增大近似线性下降，$\eta$ 每增加 $5\%$，$S_{008}$ 上的 $q^{\max}$ 约下降 $4\!\sim\!5\,$kg。当 $\eta$ 升至约 $30\%$ 时，C 型在最远的 $S_{008}$ 处 $q^{\max}$ 逼近其满载临界，可能触发该服务区需拆分为更多架次。安全余量的载荷敏感性如图~\ref{fig:sens} 所示。

\begin{figure}[htbp]
  \centering
  \subfloat[$q^{\max}$ 随安全余量变化]{\includegraphics[width=0.48\textwidth]{q1_1_03.png}}
  \hfill
  \subfloat[组批架次数对安全余量的响应]{\includegraphics[width=0.48\textwidth]{q1_1_04.png}}
  \caption{返航安全余量敏感性分析}
  \label{fig:sens}
\end{figure}

对组批结果：在 $\eta\in[10\%,25\%]$ 区间内，最优组批架次数稳定为 $N=18$，仅 C 型重载批次的边际载荷有微调，方案结构不变；当 $\eta\ge30\%$ 时，个别 C 型重载批需拆箱，架次数上升至 $19\!\sim\!20$，总能耗随之抬升约 $2\!\sim\!4\%$。这表明基准安全策略下问题一结论稳健，而安全余量是运输能力的“硬边界”——它对总体架次的影响存在明显的阈值效应，$25\%$ 附近是方案结构的转折点。

\subsection{其他参数与算法稳定性}

对 NSGA-II（问题二、三）改变随机种子与种群规模（$60/80/100$）重复求解，代表解的架次数与能耗波动均在 $3\%$ 以内，非支配前沿形态一致，说明启发式求解稳定、结论可复现。对通信参数（衰落裕量 $M$、遮挡损耗 $L_{obs}$）在 $\pm2\,$dB 内扰动，需中继相位数变化不超过 $5\%$，且候选悬停网格仍能实现零通信缺口，表明问题三的中继覆盖方案对链路预算的中等不确定性具有鲁棒性。
